
\subsection{A finite sphere-only diagram with every map specified}

\begin{result}{B.30}{Proved}{A geometric word-colimit and its exact inverse}
For every pinned stable graph, its framed fundamental groupoid
has an ordinary finite-colimit presentation using only finite
ray models of three- and four-point spheres, after the explicit
central induction of B.1.  Every arrow of the diagram is one of
the concrete two-generator path maps below.  The colimit has a
free action on its sign objects, and its comparison to the
framed fundamental groupoid has the displayed inverse.

The geometric generators may be chosen from B.27--B.29, so
this presentation includes the elliptic cut changes and their
full framing relations.  The low-type diagrams produced by
\texttt{b\_sphere\_colimit.py} have the following sizes; these
are diagram blocks, rather than the objects of their colimits:
\[
\begin{array}{c|rrrr}
(g,n)&\text{three-point blocks}&\text{four-point blocks}
 &\text{multiplication blocks}&\text{input relators}\\\hline
(0,3)&2&0&0&0\\
(0,4)&2&2&0&0\\
(0,5)&2&99&44&6\\
(1,1)&2&40&18&2\\
(1,2)&2&63&28&4 .
\end{array}                                                     \tag{B.98}
\]
Multiplication blocks are included among the four-point blocks.
All generator images, terminal equations and comparison words
are recorded in the program's JSON output.

There is also an explicit three-object ordered index after
taking finite coproducts of these blocks.  The comparison with
the particular uncompleted reduced-tree assignment of the
original scan~100 remains the precise issue recorded in B.8:
the construction here specifies its own index and its own path
maps.  In genus zero, B.27 supplies the separate geometric
reduced-tree diagram, and the substitutions below compare the
two resulting groupoids on generators.
\end{result}
\begin{proof}
Write $A=\mathbb Z^r$ for the retained external framing lattice
and $E=\tfrac12 A$, with sign set $\mu=E/A$.  First specify the
two elementary blocks, including their geometric maps.  At
$(0,4)$ normalize the labelled points to $0,1,\infty,\lambda$.
The coordinate tangent frames $\partial_z$ at the finite
points and $\partial_{1/z}$ at infinity trivialize all four
tangent lines over $\lambda\in\mathbb C\setminus\{0,1\}$.
Its framed group on a full orbit of sign objects is therefore
the translation groupoid of
$\tfrac12\mathbb Z^4\times F(u,v)$.  The loops $u,v$ are the
positive based meridians of $0,1$, based at $\lambda=1/2$.
Induce the central lattice along the zero map
$\mathbb Z^4\to A$.  B.1 gives the block
\[
 \mathcal P_A=\bigl(E\times F(u,v)\rightrightarrows\mu\bigr).
                                                               \tag{B.99}
\]
The same induction from a three-point sphere gives
\[
 \mathcal Z_A=\bigl(E\rightrightarrows\mu\bigr).                \tag{B.100}
\]
The notation denotes translation by the character to $E/A$;
the free generators have character zero.  The object-set
induction sends a free sign orbit to the free orbit $\mu$.
Thus (B.99)--(B.100) are the stated elementary sphere models
with explicitly changed central lattices.

There is a concrete continuous realization of every free-group
map used below.  Take the circles of radius $1/4$ around $0$
and $1$ and join them by the real segment from $1/4$ to $3/4$.
The punctured plane has this two-loop graph as a spine: a
regular neighborhood is a three-holed sphere, and its three
complementary annular ends retract into its boundary.  Choose
the collapse of that neighborhood to its indicated spine.
A map of the spine that traverses prescribed words on its two
loops extends, by this retraction, to a continuous map of the
four-point parameter space.  On the frame torus it is the
identity.  A central factor means concatenating the indicated
integer boundary rotations.  This defines the actual
fundamental-groupoid functor.  In particular the maps
\[
 (u,v)\mapsto(t,1),\quad (t,w)\mapsto(u,1),\quad
 (t,w)\mapsto(v^{\pm1},1),\quad
 (t,w)\mapsto(uv,1)                                           \tag{B.101}
\]
are maps of these specified path models, fixing the central
half-turn paths and every sign object.

Now take the finite presentation obtained by B.27, replace its
low-type generators by B.28--B.29, and choose paths from its
finitely many objects to one reference object.  Its based
framed group $H$ has a presentation
\[
 H=\left\langle A,x_1,\ldots,x_m\ \middle|\
 A\text{ central},\quad R_j(x)=c_j\ (1\leq j\leq s)
 \right\rangle,
 \qquad c_j\in A.                                             \tag{B.102}
\]
These are actual geometric generators and relations.  To
obtain (B.102), retain the central boundary generators of
B.27 and commute them to the end of each relator.  The
exponent vector at that end is $c_j$.  Local braid,
pentagon, cut-change, graph-lift and incidence relations are
all supplied in B.27--B.29.  This procedure does not require
an unknown central lift to be inferred from an abstract
presentation.  Its presentation is exact by B.27's explicit
word inverse.  Equivalently, the lifting proof of (B.6)
verifies (B.102) directly.

Here is a finite diagram all of whose blocks are
(B.99)--(B.100); no large free group is left as an undeclared
primitive.  Use two copies $C,K$ of $\mathcal Z_A$, with
$C\to K$ the identity on $E$.  For each generator $x_i$
use a copy $P_{x_i}$ of $\mathcal P_A$, whose free generators
are named $t,w$.  Give every such block the idempotent
endomorphism
\[
       \rho:t\mapsto t,\quad w\mapsto1,\quad E\mapsto E.
                                                               \tag{B.103}
\]
It is a genuine idempotent functor, not a formal relation
on objects.  Its colimit identifies $w$ with $1$.
Include the central inclusion from $C$ to each block.

Write a nonempty relator as
$R_j=x_{i_1}^{\epsilon_1}\cdots x_{i_l}^{\epsilon_l}$.
For each $0\leq k\leq l$ add a variable block $P_{p_{jk}}$
with the same endomorphism (B.103).  Add maps
\[
 P_{p_{j0}}\longrightarrow K:\ t\mapsto1,
 \qquad
 P_{p_{jl}}\longrightarrow K:\ t\mapsto c_j,
                                                               \tag{B.104}
\]
which kill $w$ and fix $E$.  For each $1\leq k\leq l$
add a multiplication block $M_{jk}=\mathcal P_A$ with
free generators $u,v$ and the three maps
\[
\begin{array}{c|c}
\text{source}&\text{image of its generator }t\\\hline
P_{p_{j,k-1}}&u\\
P_{x_{i_k}}&v^{\epsilon_k}\\
P_{p_{jk}}&uv.
\end{array}                                                     \tag{B.105}
\]
They kill $w$ and fix $E$.  Include $C\to M_{jk}$.
The index has just these arrows, identities, the idempotents,
and their forced composites.  Impose $f\rho=f$ for the
outgoing maps and identify all composites of central
inclusions with the stated central inclusion.  It is finite:
its only nonidentity loops are the idempotents, and after
leaving a variable block every remaining arrow ends at a
multiplication block or at $K$.  These composition rules
hold for the actual functors (B.103)--(B.105).
Empty relators may be omitted: by the injection of $A$ into
$H$, an empty relator has $c_j=0$.

The colimit can now be computed without an abstract existence
argument.  Relation (B.103) removes every dummy $w$.
Equation (B.105) gives, in the colimit,
\[
 u=p_{j,k-1},\qquad v=x_{i_k}^{\epsilon_k},\qquad
 p_{jk}=p_{j,k-1}x_{i_k}^{\epsilon_k}.                          \tag{B.106}
\]
The first equation in (B.104) starts this recursion at $1$.
The second ends it at $c_j$.  Eliminating the $p_{jk},u,v$
therefore gives exactly $R_j=c_j$, with no other relation
on the $x_i$ or on $E$.  The colimit is the translation
groupoid of
\[
 H^{1/2}=(H\times E)/\{(a,-a):a\in A\}
          \quad\text{on }\mu.                                \tag{B.107}
\]
All object maps throughout the diagram are the identity of
$\mu$, so its colimit has exactly these sign objects, and
the sign action remains free.

Here are the comparison functors and their inverse on every
block.  Send $t$ in $P_{x_i}$ to the geometric loop $x_i$;
send $t$ in $P_{p_{jk}}$ to the length-$k$ prefix of $R_j$;
send $u,v$ in $M_{jk}$ to the preceding prefix and to
$x_{i_k}^{\epsilon_k}$.  Send all central half turns to the
actual half turns.  The terminal maps commute because their
precise equation is $R_j=c_j$, whose geometric homotopy is
one of the supplied relator homotopies.  This gives a cone
and hence the forward functor.  In the other direction send
$x_i$ to $t$ in $P_{x_i}$ and each central element to its
copy in $C$.  Equations (B.104)--(B.106) prove each relator,
so this is an actual inverse functor on (B.102).  The two
composites are the identity on the geometric generators,
and on every auxiliary generator by (B.103) and (B.106).
Additional finite basepoints are restored using the chosen
connecting paths; those paths give the inverse natural
isomorphism.  This verifies the comparison rather than
appealing to invariance of arbitrary ordinary colimits
under equivalence.

The path substitutions in B.28--B.29 also give the promised
geometric dictionary for (B.102).  A parenthesized sphere
half turn is inserted by (B.85); an elliptic cut change is
$aba$; the cubic boundary is $(ab)^6$; and the quartic
boundary product is $(abc)^4=z_1z_2$.  The flag maps are
(B.97) and B.27's transporter conjugations.  Reading a
cell path gives the forward word, and B.27's inverse word
reading followed by these substitutions gives the inverse.
All the relator homotopies can be recovered as finite
products of those listed elementary homotopies.  An
explicit terminating procedure enumerates products of
conjugates of the finitely many defining relators until
the required null word is reached; termination follows
from the already proved presentation (B.102).  Thus this
is a construction of the particular maps and homotopies,
not an unspecified factorization through sphere groups.

For (B.98), use the five-generator six-relator presentation
B.19 for $(0,5)$, the two-generator braid and boundary
presentation (B.10) for $(1,1)$, and (B.11) for $(1,2)$.
Their total relator lengths are respectively $44,18,28$.
The variable-block counts are $55,22,35$, giving the
four-point-block totals in the table.  In genus zero the
normalization $0,1,\infty$ trivializes the tangent lines,
so the six five-point central vectors are zero.  For
$(1,1)$ the boundary relator has vector $1$; for $(1,2)$
the full-twist relator has vector $(1,1)$.  The program
verifies all substitutions, multiplication maps and
terminal equations, including these nonzero framing
vectors, by the faithful braid calculations.

Finally the asserted ordered index is completely explicit.
For the finite diagram $D:\mathcal I\to\mathrm{Gpd}$ just
constructed, put
\[
 U=\coprod_{\alpha:i\to j}D(i),\qquad
 V=\coprod_{i\in\mathcal I}D(i).
\]
Let $f:U\to V$ be the source inclusion and $g:U\to V$ the
functor $D(\alpha)$ followed by target inclusion on the
$\alpha$ summand.  Its colimit is the coequalizer of $f,g$:
a cocone is precisely a map from $V$ agreeing on those two
images of every arrow.  This coequalizer is the pushout
\[
       U\ \xleftarrow{\mathrm{fold}}\ U\amalg U
           \ \xrightarrow{\ f\amalg g\ }\ V.                 \tag{B.108}
\]
Indeed a pushout cocone requires its map on $U$ to equal
both composites with $f,g$, which is exactly the same
universal condition.  Thus (B.108) is the advertised
three-object poset diagram, with its actual functors.
Its values are finite coproducts of the elementary sphere
blocks and its colimit is still (B.107).  This ordered
index and B.27's reduced-tree index are related by their
explicit comparison and inverse to (B.107).  An
identification of their individual values with the
unstated original tree assignment is not supplied by
renaming either index; it is the remaining source
comparison in B.8.
\end{proof}
